\lyseis{1.5}{Χαρακτήρες διαιρετότητας, πρώτοι αριθμοί, ΜΚΔ και ΕΚΠ}
\sol{1} Με το 2: 2.480, 6.002, 1.050. Με το 5: 2.480, 3.115, 1.050. Με το 10: 2.480, 1.050.
\sol{2} Άθροισμα ψηφίων: 252 $\to$ 9 (ναι), 3.614 $\to$ 14 (όχι), 8.001 $\to$ 9 (ναι), 55.551 $\to$ 21 (ναι), 1.370 $\to$ 11 (όχι). Έλεγχος, π.χ. $252:3=84$.
\sol{3} 738 $\to$ 18 (ναι), 5.427 $\to$ 18 (ναι), 92.808 $\to$ 27 (ναι), 6.661 $\to$ 19 (όχι), 14.445 $\to$ 18 (ναι).
\sol{4} 524 (24, ναι), 1.018 (18, όχι), 6.300 (00, ναι), 3.772 (72, ναι), 90 (το 90 δεν διαιρείται με το 4, όχι).
\sol{5} 450 (50, ναι), 2.125 (25, ναι), 3.300 (00, ναι), 1.080 (80, όχι), 7.675 (75, ναι).
\sol{6} \begin{tabular}[t]{@{}lccccccc@{}}\toprule & 2 & 3 & 4 & 5 & 9 & 10 & 25\\\midrule 360 & ναι & ναι & ναι & ναι & ναι & ναι & όχι\\ 1.125 & όχι & ναι & όχι & ναι & ναι & όχι & ναι\\ 5.724 & ναι & ναι & ναι & όχι & ναι & όχι & όχι\\ 40.500 & ναι & ναι & ναι & ναι & ναι & ναι & ναι\\\bottomrule\end{tabular}
\sol{7} Το 78 (άρτιος και $7+8=15$), το 3.402 (άρτιος και $3+4+0+2=9$) και το 9.996 (άρτιος και $9+9+9+6=33$). Το 250 δεν διαιρείται με το 3 ($2+5+0=7$), και το 1.113 είναι περιττός.
\sol{8} α) $6+3+8=17$, χρειάζεται 18, άρα $\square=1$: 6.318. \quad β) $4+2=6$ και πρέπει να φτάσει σε πολλαπλάσιο του 3: $\square=0,\,3,\,6,\,9$ (402, 432, 462, 492). \quad γ) Τα δύο τελευταία ψηφία $\square 6$ πρέπει να διαιρούνται με το 4: $16, 36, 56, 76, 96$, άρα $\square=1,\,3,\,5,\,7,\,9$. \quad δ) Πρέπει να τελειώνει σε 0 (διαιρείται με 2 και 5): 290.
\sol{9} α) Λάθος: το 16 τελειώνει σε 6 και δεν διαιρείται με το 6. β) Σωστό: το 9 διαιρείται με το 3. γ) Σωστό. δ) Λάθος: το 24 διαιρείται και με τα δύο, αλλά όχι με το 48. ε) Σωστό. στ) Σωστό.
\sol{10} Το 4.752 έχει άθροισμα ψηφίων 18, άρα διαιρείται με το 3 και το 9. Τελειώνει σε 2, άρα δεν διαιρείται με το 5. Τα δύο τελευταία ψηφία 52 διαιρούνται με το 4. Είναι άρτιος και διαιρείται με το 3, άρα και με το 6. Σε 3 ομάδες: $4.752:3=1.584$ €. Σε 4: $1.188$ €. Σε 6: $792$ €. Σε 9: $528$ €. Σε 5 όχι.
\sol{11} Το 150 τελειώνει σε 0, άρα διαιρείται με το 2, το 5 και το 10. Το άθροισμα ψηφίων είναι 6, άρα διαιρείται με το 3, και επειδή είναι και άρτιος, με το 6. Τα δύο τελευταία ψηφία 50 δεν διαιρούνται με το 4. Γίνεται με 2, 3, 5, 6 και 10 στοίβες.
\sol{12} Ο αριθμός είναι $54\square 0$, με άθροισμα ψηφίων $9+\square$. Πρέπει να διαιρείται με το 9, άρα $\square=0$ ή $9$. Επειδή το τρίτο ψηφίο είναι μεγαλύτερο από το 4, είναι το 9: ο κωδικός είναι 5.490.
\sol{13} Με το 5: $90:5=18$ θέσεις. Με το 10: 9 θέσεις (10, 20, \dots, 90). Και με το 3 και με το 5: οι 15, 30, 45, 60, 75, 90, δηλαδή 6 θέσεις.
\sol{14} Το 2.024 είναι δίσεκτο ($24:4=6$). Το 2.000 είναι δίσεκτο (διαιρείται με το 400). Το 1.900 δεν είναι (διαιρείται με το 100 αλλά όχι με το 400). Το 2.100 δεν είναι. Το 2.030 δεν είναι (30 δεν διαιρείται με το 4).
\sol{15} Αν ένας αριθμός τελειώνει σε 0, είναι πολλαπλάσιο του 10: $10\cdot\kappa$. Επειδή $10=2\cdot 5$, έχουμε $10\kappa=2\cdot(5\kappa)=5\cdot(2\kappa)$. Άρα διαιρείται και με το 2 και με το 5.
\sol{16} Τα κομμάτια $2\cdot 999$, $7\cdot 99$ και $6\cdot 9$ διαιρούνται με το 9 και με το 3. Αρκεί να ελέγξουμε το $2+7+6+3=18$. Το 18 διαιρείται και με το 3 και με το 9, άρα ο 2.763 διαιρείται και με τα δύο ($2.763=9\cdot 307$).
\sol{17} Πρέπει να διαιρείται με το 4, το 5 και το 6. Τα πολλαπλάσια του $\EKP(4,\,5,\,6)=60$ είναι 60, 120, \dots, και το πρώτο τετραψήφιο είναι $60\cdot 17=1.020$.
\sol{18} Το μεσαίο ψηφίο είναι το άθροισμα των δύο άκρων: $1+1=2$, $1+2=3$, $1+3=4$, $1+4=5$. Το 165 ταιριάζει ($1+5=6$) και $165=11\cdot 15$. Το 175 δεν ταιριάζει ($1+5=6\neq 7$) και $175:11$ δίνει υπόλοιπο 10.
\sol{19} Για το 4 δεν αρκεί το τελευταίο ψηφίο: κοιτάμε τα δύο τελευταία ψηφία μαζί. Το 34 δεν διαιρείται με το 4 ($34=4\cdot 8+2$), άρα ο 5.034 δεν διαιρείται με το 4.
\sol{20} 16: 1, 2, 4, 8, 16. \quad 20: 1, 2, 4, 5, 10, 20. \quad 28: 1, 2, 4, 7, 14, 28. \quad 42: 1, 2, 3, 6, 7, 14, 21, 42. \quad 49: 1, 7, 49.
\sol{21} Του 3: 0, 3, 6, 9, 12, 15. Του 8: 0, 8, 16, 24, 32, 40. Του 15: 0, 15, 30, 45, 60, 75.
\sol{22} Πρώτοι: 3, 23, 31, 43, 59, 83. Το 1 δεν είναι ούτε πρώτος ούτε σύνθετος. Σύνθετοι: 15 (3), 27 (3), 39 (3), 49 (7), 51 (3), 63 (3), 77 (7).
\sol{23} 83, 89, 97, 101, 103, 107, 109.
\sol{24} α) $36=2^2\cdot 3^2$ \quad β) $90=2\cdot 3^2\cdot 5$ \quad γ) $126=2\cdot 3^2\cdot 7$ \quad δ) $200=2^3\cdot 5^2$ \quad ε) $360=2^3\cdot 3^2\cdot 5$ \quad στ) $462=2\cdot 3\cdot 7\cdot 11$.
\sol{25} α) $16\cdot 3=48$ \quad β) $2\cdot 25\cdot 7=350$ \quad γ) $9\cdot 5\cdot 7=315$ \quad δ) $8\cdot 3\cdot 11=264$.
\sol{26} α) Σωστό. β) Λάθος. γ) Λάθος: $77=7\cdot 11$. δ) Λάθος: $3+5=8$. ε) Σωστό: έχει διαιρέτες τουλάχιστον τους δύο παράγοντες. στ) Σωστό: $0=\alpha\cdot 0$. ζ) Λάθος: $57=3\cdot 19$.
\sol{27} Ο 323 δεν διαιρείται με το 2, το 3 ($3+2+3=8$), το 5, το 7 ($7\cdot 46=322$), το 11 ($11\cdot 29=319$) ούτε το 13 ($13\cdot 24=312$). Με το 17: $17\cdot 19=323$. Οι πρώτοι είναι 17 και 19.
\sol{28} Οι διαιρέτες του 54 είναι 1, 2, 3, 6, 9, 18, 27, 54: σακουλάκια με όποιον από αυτούς τους αριθμούς καραμελών. Το 53 είναι πρώτος, άρα έχει μόνο τους διαιρέτες 1 και 53: δύο μόνο τρόποι.
\sol{29} Οι διαιρέτες του 72 είναι 1, 2, 3, 4, 6, 8, 9, 12, 18, 24, 36, 72. Αφαιρώ τον 1 και τον 72: μένουν 10, άρα 10 τρόποι (2, 3, 4, 6, 8, 9, 12, 18, 24 ή 36 παιδιά σε κάθε ομάδα). Το 73 είναι πρώτος, άρα κανένας τρόπος.
\sol{30} Οι πρώτοι ανάμεσα στο 30 και στο 50 είναι οι 31, 37, 41, 43 και 47. Άθροισμα ψηφίων: 4, 10, 5, 7, 11. Μόνο ο 43 έχει άθροισμα 7: ο αριθμός είναι ο 43.
\sol{31} Οι πρώτοι μέχρι το 60 είναι 17: 2, 3, 5, 7, 11, 13, 17, 19, 23, 29, 31, 37, 41, 43, 47, 53, 59. Το 2 σβήνει 29 αριθμούς (4, 6, 8, \dots, 60). Το 3 σβήνει ακόμα 9 αριθμούς: 9, 15, 21, 27, 33, 39, 45, 51, 57 (τα άλλα πολλαπλάσιά του είναι ήδη σβησμένα, γιατί είναι άρτια).
\sol{32} Για το 6: $1+2+3=6$. Για το 28: $1+2+4+7+14=28$. Σε κάθε περίπτωση το άθροισμα των διαιρετών (εκτός από τον ίδιο) ισούται με τον αριθμό.
\sol{33} $(3,\,5)$, $(5,\,7)$, $(11,\,13)$, $(17,\,19)$, $(29,\,31)$, $(41,\,43)$.
\sol{34} Οι αριθμοί είναι $1, 3, 9, 2, 6, 18, 4, 12, 36, 8, 24, 72$, δηλαδή 12 αριθμοί. Είναι όλοι οι διαιρέτες του 72.
\sol{35} Αν το 1 ήταν πρώτος, το 6 θα είχε πολλές αναλύσεις: $2\cdot 3$, $1\cdot 2\cdot 3$, $1\cdot 1\cdot 2\cdot 3$, \dots, όσες θέλεις. Η ανάλυση δεν θα ήταν μοναδική.
\sol{36} $101\cdot 103=10.100+303=10.403$. Αν δοκιμάζαμε τους πρώτους με τη σειρά (2, 3, 5, \dots, 97), θα αποτυγχάναμε 25 φορές, και μόνο ο 101 θα πετύχαινε. Με ακόμα μεγαλύτερους αριθμούς η δουλειά γίνεται τεράστια, και σε αυτό στηρίζεται η ασφάλεια των κωδικών.
\sol{37} α) Διαιρέτες του 14: 1, 2, 7, 14. Του 35: 1, 5, 7, 35. $\MKD=7$. \quad β) Του 24: 1, 2, 3, 4, 6, 8, 12, 24. Του 36: 1, 2, 3, 4, 6, 9, 12, 18, 36. $\MKD=12$. \quad γ) Του 16: 1, 2, 4, 8, 16. Του 45: 1, 3, 5, 9, 15, 45. Κοινός μόνο ο 1: $\MKD=1$. \quad δ) Του 28: 1, 2, 4, 7, 14, 28. Του 42: 1, 2, 3, 6, 7, 14, 21, 42. $\MKD=14$.
\sol{38} α) $\EKP(6,\,8)=24$ \quad β) $\EKP(9,\,12)=36$ \quad γ) $\EKP(10,\,15)=30$ \quad δ) $\EKP(4,\,9)=36$.
\sol{39} α) $30=2\cdot 3\cdot 5$, $45=3^2\cdot 5$: $\MKD=3\cdot 5=15$, $\EKP=2\cdot 3^2\cdot 5=90$. \quad β) $72=2^3\cdot 3^2$, $108=2^2\cdot 3^3$: $\MKD=2^2\cdot 3^2=36$, $\EKP=2^3\cdot 3^3=216$. \quad γ) $84=2^2\cdot 3\cdot 7$, $126=2\cdot 3^2\cdot 7$: $\MKD=2\cdot 3\cdot 7=42$, $\EKP=2^2\cdot 3^2\cdot 7=252$. \quad δ) $150=2\cdot 3\cdot 5^2$, $225=3^2\cdot 5^2$: $\MKD=3\cdot 5^2=75$, $\EKP=2\cdot 3^2\cdot 5^2=450$.
\sol{40} $20=2^2\cdot 5$, $30=2\cdot 3\cdot 5$, $50=2\cdot 5^2$. $\MKD=2\cdot 5=10$ (κοινοί παράγοντες 2 και 5, με μικρότερο εκθέτη 1). $\EKP=2^2\cdot 3\cdot 5^2=300$.
\sol{41} Πρώτοι μεταξύ τους είναι τα (α), (β), (δ) και (στ). Τα (γ) 18 και 27 έχουν κοινό διαιρέτη το 9, και τα (ε) 35 και 49 το 7.
\sol{42} α) $\MKD=6$, $\EKP=48$ \quad β) $13$ και $52$ \quad γ) $20$ και $100$. Όταν ο ένας αριθμός διαιρεί τον άλλον, το ΜΚΔ είναι ο μικρότερος και το ΕΚΠ ο μεγαλύτερος.
\sol{43} α) Λάθος: το ΜΚΔ είναι 4. β) Λάθος: το ΕΚΠ είναι 18. γ) Σωστό. δ) Σωστό. ε) Σωστό. στ) Σωστό.
\sol{44} $10=2\cdot 5$ και $25=5^2$, άρα $\EKP=2\cdot 5^2=50$ δευτερόλεπτα. Στα 3 λεπτά (180 δευτ.) συμπίπτουν στα 50, 100 και 150 δευτερόλεπτα: 3 φορές.
\sol{45} $90=2\cdot 3^2\cdot 5$, $150=2\cdot 3\cdot 5^2$, άρα $\MKD=2\cdot 3\cdot 5=30$. Το πολύ 30 καλάθια, με $90:30=3$ μήλα και $150:30=5$ πορτοκάλια το καθένα.
\sol{46} $240=2^4\cdot 3\cdot 5$, $180=2^2\cdot 3^2\cdot 5$, άρα $\MKD=2^2\cdot 3\cdot 5=60$. Μέγιστη πλευρά 60 εκ. Πλακάκια: $240:60=4$ στο μήκος και $180:60=3$ στο πλάτος, δηλαδή $4\cdot 3=12$ πλακάκια.
\sol{47} $12=2^2\cdot 3$ και $16=2^4$, άρα $\EKP=2^4\cdot 3=48$ λεπτά: ξανά μαζί στις 7:48. Μαζί βρίσκονται στις 7:00, 7:48, 8:36 και 9:24 (το επόμενο θα ήταν 10:12, μετά τις 10:00): 4 φορές.
\sol{48} $\EKP(4,\,6,\,9)$: $4=2^2$, $6=2\cdot 3$, $9=3^2$, άρα $2^2\cdot 3^2=36$. Σε 36 μέρες.
\sol{49} $(10,\,25)$: ΜΚΔ 5, ΕΚΠ 50, $5\cdot 50=250=10\cdot 25$. $(16,\,24)$: 8 και 48, $8\cdot 48=384=16\cdot 24$. $(21,\,35)$: 7 και 105, $7\cdot 105=735=21\cdot 35$. $(18,\,30)$: 6 και 90, $6\cdot 90=540=18\cdot 30$. Το γινόμενο ΜΚΔ·ΕΚΠ είναι πάντα ίσο με το γινόμενο των δύο αριθμών.
\sol{50} $180=84\cdot 2+12$, και $\MKD(180,\,84)=12=\MKD(84,\,12)$. Συνεχίζω: $84=12\cdot 7+0$. Ο τελευταίος διαιρέτης πριν το υπόλοιπο 0, το 12, είναι το ΜΚΔ.
\sol{51} $9=3^2$ και $20=2^2\cdot 5$ δεν έχουν κοινό παράγοντα. Το ΕΚΠ πρέπει να περιέχει όλους τους παράγοντες και των δύο, με τον μεγαλύτερο εκθέτη, και επειδή δεν υπάρχει κανένας κοινός, τους παίρνει όλους: $2^2\cdot 3^2\cdot 5=180=9\cdot 20$.
\sol{52} Το γινόμενο των δύο αριθμών πρέπει να είναι $4\cdot 48=192$. Ένα ζευγάρι είναι το 4 και το 48 (ΜΚΔ 4, ΕΚΠ 48). Δεύτερο ζευγάρι: 12 και 16 (ΜΚΔ 4, ΕΚΠ 48, γινόμενο 192).
\sol{53} Όχι. Το ΜΚΔ είναι μικρότερο ή ίσο από τον μικρότερο αριθμό και το ΕΚΠ μεγαλύτερο ή ίσο από τον μεγαλύτερο. Αν οι αριθμοί είναι διαφορετικοί, ο μικρότερος είναι αυστηρά μικρότερος από τον μεγαλύτερο, άρα $\MKD<\EKP$. Ισότητα έχουμε μόνο όταν οι δύο αριθμοί είναι ίσοι.
\sol{54} $\EKP(5,\,6,\,8)$: $5$, $6=2\cdot 3$, $8=2^3$, άρα $2^3\cdot 3\cdot 5=120$. Σε 120 ημέρες.
\sol{55} Το πλήθος είναι πολλαπλάσιο του 4, του 6 και του 9, άρα του $\EKP(4,\,6,\,9)=36$. Το λιγότερο έχει 36 κουλουράκια. Τα πολλαπλάσια του 36 είναι 36, 72, 108, 144, 180, 216, 252, \dots, άρα ανάμεσα στο 150 και στο 250 είναι 180 ή 216.
\sol{56} $\EKP(6,\,8)=24$ ημέρες. Η Αγγελική πάει $24:6=4$ φορές μέχρι τότε, άρα 3 φορές ανάμεσα. Ο Νίκος $24:8=3$ φορές, άρα 2 φορές ανάμεσα.
\sol{57} Οι κοινοί διαιρέτες δύο αριθμών είναι ακριβώς οι διαιρέτες του ΜΚΔ τους. Ο 18 έχει διαιρέτες και τους 2, 3, 6, 9, άρα και οι δύο αριθμοί διαιρούνται με όλους αυτούς.
\sol{58} Σύνθετος. Το τριπλάσιο $3p$ έχει διαιρέτες τουλάχιστον τους $1$, $3$ και $3p$, και ο 3 είναι διαφορετικός από το 1 και από το $3p$ (αφού $3p\geq 6$). Άρα έχει διαιρέτη εκτός από το 1 και τον εαυτό του.
\sol{59} α) 1, 2, 4, 5, 8, 10, 20, 40 \quad β) 1, 2, 4, 7, 8, 14, 28, 56 \quad γ) 1, 2, 3, 6, 11, 22, 33, 66 \quad δ) 1, 7, 13, 91 \quad ε) 1, 2, 3, 5, 6, 10, 15, 25, 30, 50, 75, 150 \quad στ) 1, 13, 17, 221.
\sol{60} α) $66=2\cdot 3\cdot 11$ \quad β) $275=5^2\cdot 11$ \quad γ) $840=2^3\cdot 3\cdot 5\cdot 7$ \quad δ) $1.001=7\cdot 11\cdot 13$.
\sol{61} Με το 6: $42+18=60=6\cdot 10$ και $42-18=24=6\cdot 4$. Γενικά αν $\delta\mid\beta$ και $\delta\mid\gamma$, τότε $\beta=\delta\kappa$ και $\gamma=\delta\lambda$, άρα $\beta+\gamma=\delta(\kappa+\lambda)$. Για το $5\alpha+15$: ο 5 διαιρεί το $5\alpha$ και το 15, άρα και το άθροισμά τους. Ή $5\alpha+15=5(\alpha+3)$.
\sol{62} $24\cdot 101=2.424$ και $58\cdot 101=5.858$: το διψήφιο γράφεται δύο φορές. Εξήγηση: $24\cdot 101=24\cdot(100+1)=2.400+24=2.424$. Το $24\cdot 100$ βάζει το 24 στις εκατοντάδες και το $+24$ στις μονάδες.
